% Recovery provenance: reconstructed on 2026-11-20 from all three complete
% authored mathematical-text chunks sent by higher_comparison2.
% The historical 17469-character TeX byte stream was lost during the outage.
% The optional sphere-to-cube paragraph was omitted as ordered in the final audit.
\Needspace{10\baselineskip}
\subsection{A.152. Relative Hausdorff regularization of compact definable families}

\begin{saresult}{A.152}{Proved under stated hypotheses}{Compact graph families admit Lipschitz regularization by a relative homotopy}
Work over the real field in any one fixed o-minimal expansion.
All compact-set families below have bounded definable incidence sets,
nonempty compact fibres, and are continuous in Hausdorff distance.
Use maximum norms on finite-dimensional coordinate spaces.
For a based PL sphere \(\Sigma^k=\partial[-1,1]^{k+1}\), \(k\ge1\),
and such a family \(C_p\), there is a definable Lipschitz map
\(g:\Sigma^k\to\Sigma^k\), based definably homotopic to the identity,
such that \(C_{g(p)}\) is Hausdorff-Lipschitz.
For \(P=[-1,1]^k\), if \(C|_{\partial P}\) is Hausdorff-Lipschitz,
there is a definable Lipschitz \(g:P\to P\), equal to the identity on
\(\partial P\), such that \(C_g\) is Hausdorff-Lipschitz.
The convex homotopy \((1-t)\operatorname{id}+tg\) is relative to
\(\partial P\).

Consequently any based jointly definable homeomorphism cube \(U\)
on a compact definable \(F\) has a based definable homotopy to
\(V_p=U_{g(p)}\) with Hausdorff-Lipschitz graphs, preserving all
spatial conditions of the original family.
No injectivity or cell-likeness of \(g\), and no homeomorphic
graph matching, is asserted.
\end{saresult}

\begin{proof}
\textbf{A compact Lipschitz resolution.}
The Main Theorem, p.~1, of the inspected 5 September 2012 manuscript
of Kocel-Cynk--Pawlucki--Valette,
\emph{A short geometric proof that Hausdorff limits are definable
in any o-minimal structure}, later Advances in Geometry 14 (2014),
49--58, DOI \(10.1515/advgeom-2013-0028\), gives a bounded definable
\(S\) and a definable Lipschitz bijection \(\phi:S\to P\) such that
\(C_\phi\) is Hausdorff-Lipschitz. Put \(K=\overline S\).
Since \(P\) is compact, \(\phi\) extends uniquely to a Lipschitz
\(\bar\phi:K\to P\); its graph closure is definable.
It is already onto because its restriction to \(S\) is onto.
The pulled family also extends, and continuity of the original
family identifies this extension with \(C_{\bar\phi}\).
Write \(L_\phi,L_C\) for the resulting Lipschitz constants.
No manifold, ball, or triangulable-map assertion about \(K\) or
\(\bar\phi\) is used.

\textbf{The actual Euclidean metric on a finite join.}
Put \(m=k+1\) and
\[
 T_m=\{(s_1,\ldots,s_m)\in K^m:
                  \bar\phi(s_1)=\cdots=\bar\phi(s_m)\}.
\]
Let \(J\) be the image of \(T_m\times\Delta^{m-1}\) under
\[
 (s,t)\longmapsto(t_i,t_i s_i)_{i=1}^m.
 \tag{152.1}
\]
This is a compact Euclidean definable set. Inactive slots forget
their \(s_i\); conversely fill them by an active point.
Thus
\[
 \pi(z)=\bar\phi(s_i),\qquad
 \mathcal C(z)=C_{\bar\phi(s_i)}
 \tag{152.2}
\]
for any active slot are well-defined and continuous by the compact
quotient construction. The projection and incidence family are
definable. The fibre \(\pi^{-1}(p)\) is precisely the labelled
join \(K_p^{*m}\).

Let \(M=\max_K\|s\|_\infty\), and let \(D_P,D_C\) bound respectively
the diameter of \(P\) and all Hausdorff fibre distances.
For \(z=(t_i,v_i)\), \(z'=(t'_i,v'_i)\), write
\(\delta=\|z-z'\|_\infty\). If \(\delta<1/(2m)\), choose
\(t_i\ge1/m\). Then \(t_i'>0\), and
\[
 \left\|\frac{v_i}{t_i}-\frac{v_i'}{t_i'}\right\|_\infty
 \le
 \frac{\|v_i-v_i'\|_\infty}{t_i}
 +\frac{\|v_i'\|_\infty|t_i-t_i'|}{t_it_i'}
 \le m(1+M)\delta.
 \tag{152.3}
\]
For larger \(\delta\), use the diameter bounds.
Consequently \(\pi\) and \(\mathcal C\) are globally Lipschitz,
with respective constants
\[
 \max\{m(1+M)L_\phi,2mD_P\},
 \qquad
 \max\{m(1+M)L_C,2mD_C\}.
 \tag{152.4}
\]
This controls the actual weighted-coordinate metric on the join.

\textbf{Connectivity and the proper mapping theorem.}
Each \(K_p\) is a nonempty compact definable set, hence a finite CW
space. Two-fold joins are path connected; joins of at least three
nonempty spaces are simply connected, by van Kampen.
Choosing a base vertex in each factor, removal of the contractible
basepoint cones identifies the join up to homotopy with the
appropriate suspended smash product. Accordingly its reduced
cellular chains are chain homotopy equivalent to the
\((m-1)\)-shifted tensor product of the reduced based complexes.
Those complexes are concentrated in nonnegative degrees, so reduced
homology vanishes below degree \(m-1\).
Hurewicz induction, together with the preceding connectivity,
shows that \(K_p^{*m}\) is \((m-2)=(k-1)\)-connected, including
disconnected initial fibres.

The fibres and \(J\) are locally contractible.
The set \(J\) is connected: a nontrivial compact clopen separation
would, because its fibres are connected, project to disjoint closed
sets separating the connected parameter space. It is therefore path
connected. Smale's Main Theorem, printed p.~604, and its based
surjectivity formulation, Theorem 9, p.~608, in
\emph{A Vietoris mapping theorem for homotopy},
Proceedings of the AMS 8 (1957), 604--610, apply.
The spaces are locally compact separable metric spaces, \(J\) is
locally \(k\)-connected, the fibres are locally
\((k-1)\)-connected and \((k-1)\)-connected, and \(\pi\) is proper
and onto. Thus
\[
 \pi_k(J,j_0)\longrightarrow\pi_k(P,p_0)
 \quad\hbox{is onto}
 \tag{152.5}
\]
and the maps in lower degrees are isomorphisms.
For \(k=1\) only connected two-fold join fibres are needed.
No fibration, continuous section, or triangulation of \(\pi\) is
asserted.

\textbf{The sphere case.}
Choose \(j_0\) above the basepoint \(p_*\).
Equation (152.5) gives a continuous based
\(q:\Sigma^k\to J\) such that \(\pi q\) is based ordinarily
homotopic to the identity.
Czapla, \emph{Definable triangulations with regularity conditions},
Geometry and Topology 16 (2012), 2067--2095,
Theorem 3.12(a), p.~2082, and Remark 3.14, p.~2084, give a definable
globally Lipschitz triangulation \(\tau:|L|\to J\), compatible with
\(\{j_0\}\). The preimage of \(j_0\) is a vertex.
Relative simplicial approximation of \(\tau^{-1}q\) gives a based
PL map \(a\) and an ordinary based homotopy.
Then \(q'=\tau a\) and \(g=\pi q'\) are definable Lipschitz,
\(g(p_*)=p_*\), and \(\mathcal Cq'=C_g\) is Hausdorff-Lipschitz.

Here the ordinary homotopy can be replaced explicitly by a definable
one. Let \(H:\Sigma^k\times I\to\Sigma^k\) be an ordinary based
homotopy from \(g\) to the identity. Approximate it, as a map into
\(R^{k+1}\), uniformly within \(1/4\), by a finite PL map \(A\),
relative to \(\{p_*\}\times I\), where \(A=p_*\).
On \(1/2\le\|x\|_\infty\le3/2\), put \(R(x)=x/\|x\|_\infty\).
This is semialgebraic and Lipschitz. Concatenate
\[
 R((1-s)g(p)+sA(p,0)),\qquad
 R(A(p,t)),\qquad
 R((1-s)A(p,1)+sp).
 \tag{152.6}
\]
Every argument stays in the shell by the approximation bound.
The basepoint stays fixed. These formulas are definable and jointly
Lipschitz, since \(g,\operatorname{id},A\) are Lipschitz.
No comparison theorem about families of homeomorphisms is used.

\textbf{A literal boundary section for the cube.}
Before forming the join, replace \(K\) by the separated tagged union
\[
 K^+=(\{0\}\times K\times\{0\}^k)
 \ \amalg\
 (\{1\}\times\{0\}^{\dim_{\rm amb}K}\times\partial P).
 \tag{152.7}
\]
Use \(\bar\phi\) on its first part and the identity projection on
the second, with the original \(C_p\) there.
Each part has Lipschitz projection and family. Their separation
and bounded images give global Lipschitz constants.
Construct the same \(m=k+1\) join.
There is a literal PL bi-Lipschitz section
\(b:\partial P\to J\): use weight one in slot one at the tagged
boundary point \(p\), and zero in all other slots.
Its inverse on its image is a coordinate projection.

Since \(P\) is contractible, Smale's theorem makes \(J\)
\((k-1)\)-connected. Therefore \(b\) extends continuously to a
map \(u:P\to J\); for \(k=1\), path connectedness joins the two
endpoint values.
Take a Lipschitz triangulation \(\tau\) compatible with
\(b(\partial P)\), marking the individual endpoint singletons
when \(k=1\).
Ordinary simplicial approximation as a map of pairs
\[
 (P,\partial P)\longrightarrow
 (|L|,\tau^{-1}b(\partial P))
\]
gives a PL \(a\) and a pair homotopy to \(\tau^{-1}u\).
Thus \(u'=\tau a\) is definable Lipschitz, and
\(u'|_{\partial P}=bc\), where
\(c:\partial P\to\partial P\) is definable Lipschitz and ordinarily
homotopic to the identity.
For \(k=1\), take the approximation relative to the endpoint
vertices, so its boundary is already exact.
For \(k\ge2\), the explicit radial construction (152.6) on
\(\partial[-1,1]^k\) gives a joint definable Lipschitz homotopy
\(H_\partial\) from \(c\) to the identity.
No relative Lipschitz triangulation theorem is being assumed.

\textbf{Repairing the boundary exactly.}
For \(k\ge2\), set \(\rho=\|p\|_\infty\) and define
\[
 \widehat u(p)=
 \begin{cases}
 u'(2p),&\rho\le1/2,\\
 bH_\partial(p/\rho,2\rho-1),&1/2\le\rho\le1.
 \end{cases}
 \tag{152.8}
\]
At \(\rho=1/2\) both formulas give \(bc(p/\rho)\), and on the
outer boundary the value is \(b(p)\).
Both pieces are Lipschitz; division on the annulus
\(\rho\ge1/2\) is Lipschitz. Across the interface split a joining
line at the inner cube boundary to obtain a global bound.
For \(k=1\), use \(\widehat u=u'\).
Then \(g=\pi\widehat u\) is definable Lipschitz,
equals the identity on \(\partial P\), and
\(C_g=\mathcal C\widehat u\) is Hausdorff-Lipschitz.
The convex homotopy
\((1-t)\operatorname{id}+tg\) remains in \(P\) and fixes its
boundary.

For a jointly definable homeomorphism family, take
\(C_p=\Gamma(U_p)\). Joint compact continuity gives uniform
evaluation continuity, and
\[
 d_H(\Gamma(U_p),\Gamma(U_q))
       \le\sup_x\|U_p(x)-U_q(x)\|,
\]
so the graph family is Hausdorff-continuous.
On the parameter boundary it is the constant diagonal.
The homotopy
\[
 U^t(x,p)=U(x,(1-t)p+tg(p))
\]
is jointly continuous and definable, every slice is one of the
original homeomorphisms, and all based and spatial conditions
are retained. Empty \(F\) has its unique maps.
No connectedness or positive dimension of \(F\) is required.
\end{proof}

\begin{saresult}{A.152G}{Proved under stated hypotheses}{The remaining comparison question can be restricted to Hausdorff-Lipschitz graph families}
For a compact closed PL manifold \(B\) over the real field, put \(N=B\times I\),
\(A_j=B\times\{j\}\), and let \(k\ge1\).
The following are equivalent:
\begin{enumerate}
\item \(\pi_kP_{\rm PL}(B)\to\pi_kP_{\mathcal R}(B)\) is surjective.
\item Every based definable bottom-fixed cube \(V\) whose graphs are
Hausdorff-Lipschitz is based definably homotopic to
\(W_p=s_pr_p^{-1}\), where \(r,s\) are jointly definable continuous
homeomorphism families preserving \(A_0,A_1\) setwise,
both uniformly forward parameter-Lipschitz, equal on \(A_0\times I^k\),
and equal on \(N\times\partial I^k\).
\end{enumerate}
\end{saresult}
\begin{proof}
The first condition gives a jointly PL representative, factored by
\(r=\operatorname{id}\), \(s=W\); compactness supplies its forward
bound. Conversely regularize an arbitrary cube by A.152, apply the
second condition, and then apply A.148 to its factorization.
\end{proof}

\textbf{Scope.}
A homeomorphic graph matching \(q_p:N\to\Gamma(W_p)\) would supply
the two coordinate factors with their required marking and parameter
estimates. Hausdorff control alone does not supply such a bijective
matching. A.152 establishes the unparameterized graph regularity
stage by an allowed relative homotopy; A.152G isolates the remaining
matching question. The proof here uses ordinary real compactness,
Smale's theorem, and finite PL approximation. It makes no direct
non-Archimedean claim.

